\begin{helper}
Let $G$ be a finite graph, let $\nu$ be a probability measure on
activation--priority pairs $(A,\pi)$, and fix a vertex $v$.  Assume first
that each activation atom has full measure on its priority cube:
\[
\nu\{(A,\pi): A=A_0,\ \pi(u)\in[0,1]\hbox{ for every }u\in A_0\}
=\nu\{(A,\pi):A=A_0\}
\]
for every activation set $A_0$.  Assume also that, whenever $v\in A_0$,
\[
\begin{aligned}
&\nu\{(A,\pi):A=A_0,\ \pi(u)\in[0,1]\hbox{ for every }u\in A_0,\\
&\hspace{4.5cm}
\pi(u)<\pi(v)\hbox{ for every }u\in A_0\cap N_G(v)\}\\
&=\nu\{(A,\pi):A=A_0\}\,
\operatorname{ofReal}\!\left(\int_0^1
a^{|\{u:u\in A_0,\ u\sim_G v\}|}\,da\right).
\end{aligned}
\]
Then the raw event that $v$ survives the priority comparison has measure
\[
\sum_{A_0}
{\bf 1}_{v\in A_0}\,
\nu\{(A,\pi):A=A_0\}\,
\operatorname{ofReal}\!\left(\int_0^1
a^{|\{u:u\in A_0,\ u\sim_G v\}|}\,da\right).
\]
\end{helper}
\begin{proof}
Let
\[
R_{A_0}=\{(A,\pi):A=A_0,\ \pi(u)<\pi(v)
\text{ for every }u\in A_0\cap N_G(v)\}
\]
and let $C_{A_0}$ be the same event with the additional requirement that
$\pi(u)\in[0,1]$ for every $u\in A_0$.  The event that $v$ belongs to the
algorithmic output
\[
\{u\in A:\pi(w)<\pi(u)\text{ for every }w\in A\cap N_G(u)\}
\]
is exactly the finite disjoint union of the events $R_{A_0}$ over all
activation sets $A_0$ containing $v$.  This is just unpacking membership in
the output set, using \noderef{SamplingOneVertexOutputMembership}: $v$ is
output precisely when $v\in A_0$ and the displayed priority inequalities
hold.  Distinct activation sets give disjoint events because the first
coordinate cannot equal two different finite sets at once.

It remains to replace $R_{A_0}$ by $C_{A_0}$ in measure.  The inclusion
$C_{A_0}\subseteq R_{A_0}$ is immediate.  The difference
$R_{A_0}\setminus C_{A_0}$ is contained in the part of the activation atom
$\{A=A_0\}$ lying outside the priority cube.  Since $\nu$ is a probability
measure, the activation atom has finite measure.  The full-cube hypothesis
therefore implies, by the usual finite-measure subtraction identity, that
this outside-cube part has measure zero.  Hence
\[
\nu(R_{A_0})=\nu(C_{A_0}).
\]
For each $A_0$ containing $v$, the comparison hypothesis evaluates
$\nu(C_{A_0})$ as the activation-atom measure times the integral of
$a^{|\{u:u\in A_0,\ u\sim_G v\}|}$ over $[0,1]$.  Summing these values over
the finite, disjoint partition by activation set gives exactly the asserted
formula, with the indicator ${\bf 1}_{v\in A_0}$ expressing the restriction
to activation sets containing $v$.
\end{proof}
